\par\Needspace{6\baselineskip}\medskip
\hypertarget{native-087-article068}{}\pdfbookmark[2]{68.}{article-068}
\begin{center}68.\end{center}
Die Lösung der umgekehrten Aufgabe, nämlich der Bestimmung von Länge und Breite aus Rektaszension und Deklination, beruht auf demselben sphärischen Dreieck. Die oben angegebenen Formeln lassen sich daher für diesen Zweck verwenden, indem man lediglich $b$ mit $\delta$ und $l$ mit $-\alpha$ vertauscht. Da diese Formeln häufig gebraucht werden, lohnt es sich, sie auch hier anzugeben.

Nach dem Verfahren aus Artikel 66 gilt
\[
\begin{aligned}
\sin(45^0-\tfrac12 b)\sin\tfrac12(E-l)&=\cos(45^0+\tfrac12\alpha)\sin(45^0-\tfrac12(\epsilon+\delta))\\[.5em]
\sin(45^0-\tfrac12 b)\cos\tfrac12(E-l)&=\sin(45^0+\tfrac12\alpha)\cos(45^0-\tfrac12(\epsilon-\delta))\\[.5em]
\cos(45^0-\tfrac12 b)\sin\tfrac12(E+l)&=\sin(45^0+\tfrac12\alpha)\sin(45^0-\tfrac12(\epsilon-\delta))\\[.5em]
\cos(45^0-\tfrac12 b)\cos\tfrac12(E+l)&=\cos(45^0+\tfrac12\alpha)\cos(45^0-\tfrac12(\epsilon+\delta)).
\end{aligned}
\]

Nach dem anderen Verfahren aus Artikel 67 bestimmen wir den \SourceLeaf{088}Hilfswinkel $\zeta$ aus der Gleichung
\[
\operatorname{tang}\zeta=\frac{\operatorname{tang}\delta}{\sin\alpha},
\]
\noindent und erhalten
\[
\begin{aligned}
\operatorname{tang}l&=\frac{\cos(\zeta-\epsilon)\operatorname{tang}\alpha}{\cos\zeta}\\[.5em]
\operatorname{tang}b&=\sin l\operatorname{tang}(\zeta-\epsilon).
\end{aligned}
\]
\noindent Zur Kontrolle der Rechnung kann man hinzufügen
\[
\cos b=\frac{\cos\delta\cos\alpha}{\cos l}=\frac{\cos(\zeta-\epsilon)\cos\delta\sin\alpha}{\cos\zeta\sin l}\cdot
\]
\noindent Zur Bestimmung von $E$ dienen wie im vorigen Artikel die Gleichungen
\[
\begin{aligned}
\cos E&=\frac{\sin\epsilon\cos\alpha}{\cos b}=\frac{\sin\epsilon\cos l}{\cos\delta}\\[.5em]
\cos b\cos\delta\sin E&=\cos\epsilon-\sin b\sin\delta.
\end{aligned}
\]

Die differenziellen Änderungen von $l$, $b$ ergeben sich aus den folgenden Formeln.
\[
\begin{aligned}
\mathrm{d}l&=\frac{\sin E\cos\delta}{\cos b}\,\mathrm{d}\alpha+\frac{\cos E}{\cos b}\,\mathrm{d}\delta\\[.5em]
\mathrm{d}b&=-\cos E\cos\delta\,\mathrm{d}\alpha+\sin E\,\mathrm{d}\delta.
\end{aligned}
\]
